\documentclass[12pt]{article}
\usepackage[T1]{fontenc}
\usepackage[utf8]{inputenc}
\usepackage{lmodern}
\usepackage{textcomp}
\usepackage{enumitem}
\usepackage{geometry}
\usepackage{graphicx}
\usepackage{amsmath, amssymb}
\geometry{margin=1in}
\begin{document}
\begin{center}
Sociology - Graduate Level Assessment 24 \
Total Marks: 40 \
Answer all questions.
\end{center}

\section*{Multiple Choice (10 marks)}
\begin{enumerate}[label=\arabic*.]
\item Women have been excluded from the public sphere because:
\begin{enumerate}[label=(\alph*)]
    \item industrial capitalism separated the middle class home from the workplace
    \item those who enter paid employment have been 'sidelined' into particular fields
    \item it is difficult to succeed in 'malestream' politics without compromising their femininity
    \item all of the above
\end{enumerate}
\item Comte's term 'positivism' refers to:
\begin{enumerate}[label=(\alph*)]
    \item a theory that emphasizes the positive aspects of society
    \item the precise, scientific study of observable phenomena
    \item a theory that posits difficult questions and sets out to answer them
    \item an unscientific set of laws about social progress
\end{enumerate}
\item The term inter-generational mobility refers to:
\begin{enumerate}[label=(\alph*)]
    \item movement into a different occupational category over a person's lifetime
    \item movement into different occupational categories between generations
    \item movement into a higher occupational category
    \item movement into an occupation that generates a lower income
\end{enumerate}
\item Murray thought that the 'underclass' consisted of people who:
\begin{enumerate}[label=(\alph*)]
    \item formed an inferior 'race' with low levels of intelligence
    \item lived morally unsound lives of crime and squalor
    \item were too reliant upon welfare benefits
    \item all of the above
\end{enumerate}
\item Which of these trends did the New Right not suggest as evidence of declining family values?
\begin{enumerate}[label=(\alph*)]
    \item the tendency for cohabitation before marriage
    \item the rising divorce rate
    \item the absence of fathers in many households
    \item the increasing number of single parent families
\end{enumerate}
\end{enumerate}

\section*{True / False (10 marks)}
\textit{Answer True or False}
\begin{enumerate}[label=\arabic*.]
\setcounter{enumi}{5}
\item True or False: Industrial capitalism significantly contributed to the migration of populations from rural areas to urban centers during the nineteenth century.
\item True or False: Sullivan's (2000) study indicated that men did the most housework when their wives were at home and frequently nagged them.
\item True or False: The ideology process is one of the decision-making processes identified by Domhoff in the USA.
\item True or False: Employers in 19 th century factories primarily exercised disciplinary power through the implementation of routine tasks that reduced monotony.
\item True or False: Sociology is a social science because its theories are logical, supported by empirical evidence, and subject to peer scrutiny.
\end{enumerate}

\section*{Long Form Response (20 marks)}
\textit{Answer each question with a clear and complete explanation.}
\begin{enumerate}[label=\arabic*.]
\setcounter{enumi}{10}
\item Discuss Wallerstein's concept of the 'world system' in the context of the capitalist world economy, focusing on its division of labor across diverse ethnic and cultural groups.
\item Discuss the concept of social stratification, analyzing its implications on societal structure and mobility, and how it influences individual opportunities and experiences.
\end{enumerate}

\vfill
\noindent\textit{End of Paper}
\end{document}